% Part: computability
% Chapter: recursive-functions
% Section: trees

\documentclass[../../../include/open-logic-section]{subfiles}

\begin{document}

\olfileid[fr]{cmp}{rec}{tre}
\olsection{Arbres}

Il est parfois utile de représenter les arbres par des entiers
naturels, comme nous l'avons fait pour les suites finies. Leurs
propriétés et les opérations qui portent sur eux peuvent alors
être étudiées au moyen de relations et de fonctions récursives
primitives appliquées à leurs codes. Nous utiliserons à cette
fin le codage des suites. Un arbre peut se réduire à un seul
nœud, éventuellement étiqueté, ou comporter un nœud,
éventuellement étiqueté, relié à plusieurs sous-arbres.
Ce nœud est la \emph{racine} de l'arbre; les sous-arbres qui
lui sont directement reliés sont ses \emph{sous-arbres immédiats}.

Nous codons récursivement un arbre par une suite
$\tuple{k,d_1,\dots,d_k}$, où $k$ est le nombre de
sous-arbres immédiats et où $d_1,\dots,d_k$ sont leurs
codes. Si les nœuds portent des étiquettes, nous pouvons
placer celles-ci après les codes des sous-arbres immédiats.
Un arbre réduit à un nœud d'étiquette~$l$ a ainsi pour
code $\tuple{0,l}$; un arbre dont la racine porte
l'étiquette~$l_1$ et qui possède deux sous-arbres réduits
à des nœuds d'étiquettes $l_2$ et~$l_3$ a pour code
$\tuple{2,\tuple{0,l_2},\tuple{0,l_3},l_1}$.

\begin{prop}
  \ollabel{prop:subtreeseq}
  La fonction $\fn{SubtreeSeq}(t)$, qui renvoie le code
  d'une suite contenant les codes de tous les sous-arbres
  de l'arbre codé par~$t$, est récursive primitive.
\end{prop}

\begin{proof}
  Remarquons d'abord que
  $\fn{ISubtrees}(t)=\fn{subseq}(t,1,(t)_0)$
  est récursive primitive et renvoie la suite des codes
  des sous-arbres immédiats de l'arbre~$t$.
  Définissons une fonction auxiliaire
  $\fn{hSubtreeSeq}(t,n)$ qui donne une suite de tous les
  sous-arbres situés à une distance \emph{au plus}~$n$
  de la racine. Pour $n=0$, cette suite ne contient que~$t$.
  Pour passer de~$n$ à~$n+1$, on conserve les sous-arbres
  déjà énumérés et on leur adjoint les sous-arbres immédiats
  de chacun d'eux.

  Si $f(x)$ est récursive primitive et renvoie des codes de
  suites, la fonction qui concatène les images des $k$
  premiers éléments de~$s$ est également récursive primitive :
  $g_f(s,k)=f((s)_0)\concat\dots\concat f((s)_{k-1})$,
  avec la convention $g_f(s,0)=\emptyseq$. En écrivant
  $g$ pour~$g_f$, elle se définit par récurrence primitive :
  \begin{align*}
    g(s, 0) & = \emptyseq\\
    g(s, k+1) & = g(s, k) \concat f((s)_k).
  \end{align*}
  Si $s$ est une suite de codes d'arbres,\footnote{Correction éditoriale : $k$ compte ici les éléments traités. La source les indexe de $0$ à $k$, puis emploie $k=\len{s}$, ce qui traite un élément de trop. Nous utilisons les $k$ premiers éléments avec une valeur initiale vide.}
  $h(s)=g_{\fn{ISubtrees}}(s,\len{s})$ concatène donc
  exactement les listes de leurs sous-arbres immédiats,
  y compris lorsque $s$ est vide. Nous pouvons définir
  $\fn{hSubtreeSeq}$ par
  \begin{align*}
    \fn{hSubtreeSeq}(t, 0) & = \tuple{t} \\
    \fn{hSubtreeSeq}(t, n+1) & = \fn{hSubtreeSeq}(t, n) \concat
      h(\fn{hSubtreeSeq}(t, n)).
  \end{align*}
  Le code d'un sous-arbre immédiat est strictement inférieur
  au code de l'arbre qui le contient : il figure comme
  élément de la suite qui code cet arbre, et le facteur
  premier élevé à une puissance qui représente cet élément
  est strictement supérieur à lui. Les codes
  décroissent donc strictement le long de tout chemin
  de la racine à une feuille. La distance maximale de la
  racine à une feuille est ainsi bornée par~$t$.
  Par conséquent $\fn{hSubtreeSeq}(t,t)$ fournit une suite
  contenant les codes de tous les sous-arbres de l'arbre
  codé par~$t$.\footnote{Explicitation éditoriale : la récurrence affichée conserve les listes précédentes. Elle accumule donc les distances au plus $n$, avec répétitions, contrairement au texte source qui parle seulement de la distance exacte $n$. La preuve de décroissance des codes explicite aussi la borne de profondeur.}
\end{proof}

\begin{prob}
  La définition de $\fn{hSubtreeSeq}$ donnée dans la preuve
  de la \olref[cmp][rec][tre]{prop:subtreeseq} peut répéter
  certains codes. Donner une autre définition qui garantisse
  que le code de chaque sous-arbre n'apparaisse qu'une seule
  fois dans la liste obtenue.
\end{prob}

\end{document}
